\section{Related Work and Positioning}

\subsection{ETA and Route-Aware Traffic Datasets}
The broader traffic-prediction literature has been shaped by sensor-centric and aggregate spatio-temporal datasets, as reflected in surveys of traffic prediction with big data~\cite{yu2017trafficpred}. Frameworks such as LibCity have made this landscape more reproducible by standardizing datasets, tasks, and model interfaces for urban spatial-temporal prediction~\cite{libcity2021}. However, the dominant benchmark orientation remains broad traffic forecasting rather than route-aware ETA estimation in which the planned route of each vehicle is explicit. This distinction matters because ETA tasks are inherently path-dependent, whereas many widely used resources primarily capture network state, trajectories, or aggregate flows without preserving a navigation-style route context.

\subsection{Graph-Based and Route-Aware Traffic Modeling}
Graph-based learning has become a standard approach to traffic forecasting because it aligns naturally with road-network structure. Representative models such as DCRNN and ST-GCN show the value of explicit graph formulations for capturing spatial and temporal dependencies in transportation systems~\cite{dcrnn2018,stgcn2018}. In the ETA domain, production-scale work such as ETA prediction with graph neural networks in Google Maps further illustrates the importance of structured, route-aware modeling for travel-time estimation~\cite{googlemapseta2021}. Related routing-oriented work has also explored future-traffic-aware route choice and simulative ETA estimation~\cite{voloch2021,voloch2025}. The present paper differs in emphasis: rather than proposing a new predictive architecture, it focuses on the construction of a reusable route-aware representation and the tooling required to generate and evaluate it.

\subsection{Simulation and Trajectory-Conversion Pipelines}
Microscopic traffic simulation has long provided a practical basis for controlled transportation experiments, with SUMO and TraCI offering repeatable scenarios together with external programmatic control~\cite{sumo2011,traci2018}. Such simulation environments are attractive for ETA-oriented dataset construction because they expose vehicle state, route context, and traffic evolution with high observability. Real trajectory data, by contrast, requires repair, alignment, and route inference before it can be used coherently for downstream ETA analysis. These two branches are often treated separately in practice: simulation pipelines are optimized for controllability, while trajectory pipelines are optimized for recovering structure from noisy observations. This paper positions itself at their intersection by supporting both branches in one common representation.

\subsection{Testing Environments and Paper Positioning}
Existing libraries and evaluation frameworks have improved reproducibility for traffic prediction, but they typically begin from pre-existing datasets rather than from integrated dataset construction. By contrast, the web-based component presented here is designed as a downstream evaluation layer coupled to the dataset-generation workflow. Its role is to attach a compatible predictor, register prediction outcomes, and expose structured statistical analysis, plots, and reports under controlled conditions. The paper therefore should be read neither as a new ETA-model paper nor as a generic benchmark library, but as an integrated workflow for constructing route-aware ETA-oriented datasets and validating their downstream usability.
